% ============================================================================
% BORRADOR/ESQUELETO de conclusiones — estructurado por los 4 objetivos de la
% propuesta. Contiene los resultados reales del análisis; los coautores deben
% revisar/ampliar (sobre todo la parte de implementación numérica) y pulir.
% ============================================================================

\primerparrafo Este trabajo simuló numéricamente el desarrollo de la inestabilidad de Kelvin--Helmholtz (KHI) en plasmas relativistas dentro del marco de la magnetohidrodinámica relativista resistiva (RRMHD), empleando el código \textit{Cueva}, con el fin de caracterizar el efecto de la resistividad sobre la dinámica de la inestabilidad. Las conclusiones se organizan en torno a los objetivos específicos planteados.\sebnota{+ variación de campo magnético.}\claudenota{Sí: al llegar los datos, añadir un párrafo/objetivo con la conclusión de la campaña de B (intensidad/orientación). Hasta entonces el objetivo 4 se lee como ``en curso''.}

\section{Conclusiones por objetivo}

\paragraph{(1) Efecto de la resistividad sobre la tasa de crecimiento lineal.}
Se cuantificó $\gKHI(\sigma)$ sobre un barrido de $44$ simulaciones en conductividad ($\sigma\in[100,10500]$). La transición resistivo$\to$ideal se describe con un modelo sigmoide con \emph{offset} (Sección~\ref{sec:sigmoide_v6}), cuya asíntota ideal es $C+\gamma_0=1.03\pm0.05$ ($\what=0.103$). Este valor es \emph{consistente} con la teoría lineal RMHD compresible dentro del $8\%$ (Sección~\ref{sec:teoria_v6}), con el sistema holgadamente inestable ($M_{re}=0.60<\sqrt2$). La conductividad crítica que separa los regímenes se estima en $\scrit\approx2861$, con una banda sistemática real de $\pm1944$ proveniente de tres fuentes independientes de incertidumbre (Sección~\ref{sec:scrit_v6}). \emph{[Objetivo cumplido.]}\sebnota{Ver cómo la correlación energética en $\sigma=6000$ y $\sigma=10000$, en compañía del campo magnético, habla del papel de la resistividad.}

\paragraph{(2) Efecto sobre las estructuras secundarias.}
El análisis de los diagnósticos no lineales (Secciones~\ref{sec:leyes_potencia} y~\ref{sec:globales}) muestra que la enstrofía máxima escala como $\Omaxzp\propto\sigma^{1.22}$, exponente que excede el piso de una lámina de corriente Sweet--Parker única ($\sigma^{1/2}$) y se interpreta como la firma de la reconexión resistiva con \emph{tearing} secundario. La fracción de enstrofía fuera del plano $f_{Vz}$ alcanza su máximo en la zona de transición ($\sigma\approx1400$) y se suprime en el régimen ideal, indicando que las estructuras secundarias tridimensionales emergen preferentemente cuando resistividad y advección compiten en escala comparable. \emph{[Objetivo cumplido.]}

\paragraph{(3) Efecto sobre las variables globales.}
La difusión magnética controla las variables electromagnéticas integradas (Sección~\ref{sec:globales}): la corriente máxima $J_{\max}$ crece con $\sigma$ (láminas más delgadas e intensas), el trabajo electromagnético acumulado aumenta, y la amplificación de la energía magnética propia decrece de $\sim15\%$ (resistivo) a nula (ideal, campo congelado). \emph{[Objetivo cumplido.]}\claudenota{COHERENCIA (revisión Claude): corregido ``total''$\to$``propia'' (cap.~6 aclara que el panel es la energía magnética comóvil $\tfrac12\int b^2$, no $\int B^2$ de laboratorio, que de hecho crece con $\sigma$). Mantener consistente.}\sebnota{Se puede agregar un análisis igual al del campo magnético para ver la correlación de las variables energéticas y cómo esta cambia con $\sigma$.}

\paragraph{(4) Efecto de la orientación e intensidad del campo magnético.}
Este objetivo no se aborda en el cuerpo principal de la presente monografía: requiere una campaña dedicada que varíe la intensidad ($B_0$) y la orientación de $\mathbf B$ a conductividad fija, actualmente en ejecución. Un análisis preliminar (no incluido aquí) sugiere que la tensión magnética paralela al flujo ($V_{A,\parallel}$) es el agente estabilizador dominante, pero su caracterización cuantitativa queda como continuación inmediata de este trabajo. \emph{[Objetivo pendiente --- trabajo en curso.]}\claudenota{DESACTUALIZADO Y CON FÍSICA INCORRECTA (revisión Claude, PRIORIDAD): (1) la campaña B/C YA está en el cap.~6 $\Rightarrow$ reescribir este objetivo como ``cumplido'', con sus resultados (intercambio $E_{\rm kin}\leftrightarrow E_{\rm mag}$ en anti-fase; disipación óhmica vs $\theta$/$\sigma$; supresión de $\Omega^{\rm int}_{zp}$ con la tensión). (2) ``la tensión paralela al flujo $V_{A,\parallel}$ es el agente dominante'' es FALSO según el cap.~6 (corregido): en intensidad domina la presión/magnetización ($M_{A,\parallel}$ super-Alfvénico); la tensión (vía $B_y$, Campaña B) solo importa a $\theta$ grande. Reescribir todo el párrafo.}

\section{Limitaciones del Modelo y la Simulación}
\begin{itemize}
\item \textbf{Dimensionalidad:} las simulaciones son 2D; los modos tridimensionales y la cascada turbulenta plena no se capturan.
\item \textbf{Ecuación de estado:} cierre politrópico $\Gamma=4/3$, sin composición, opacidad ni radiación.
\item \textbf{Conductividad escalar:} se ignoran efectos Hall y la anisotropía respecto al campo.
\item \textbf{Resolución de la malla:} las hojas resistivas a alta $\sigma$ se acercan al límite de la malla.
\item \textbf{Asíntota ideal:} $C+\gamma_0$ es una extrapolación (el dato de mayor $\sigma$ alcanza el $85\%$); su confirmación requiere simulaciones a $\Rmst\gtrsim500$.
\end{itemize}

\section{Trabajo Futuro}
\begin{itemize}
\item \textbf{Campaña de campo magnético (objetivo 4):}\sebnota{no es válido: ya no es ``trabajo futuro'', la campaña está en curso.}\claudenota{Correcto: mover ``Campaña de campo magnético'' de Trabajo Futuro a Resultados/Conclusiones una vez integrada (martes). Si para la entrega aún no está, dejarla como ``en curso'', no como futuro.} barrido de intensidad y orientación de $\mathbf B$ para mapear la estabilización por tensión magnética.
\item \textbf{Medición directa del plateau ideal:} $2$--$3$ simulaciones a $\Rmst\gtrsim500$ ($\sigma\sim2$--$3\times10^4$) para convertir la extrapolación de $C+\gamma_0$ en medición.
\item \textbf{Extensión dimensional:} a 3D y, eventualmente, a GRMHD.
\item \textbf{Esquemas de orden aún mayor:} (p.ej.\ cuarto orden, \parencite{mignone2024}) para reducir la disipación numérica y resolver las hojas resistivas.
\end{itemize}
